% Source: https://artofproblemsolving.com/community/c3h1925394_rational_roots_test
% Stable identity: aops-topic-1925394-q01
\textbf{(a)} By the rational root theorem, a rational root of a monic polynomial with integer coefficients and constant term $1$ must be $\pm1$. But $p(1)=3$ and $p(-1)=-1$, so there is no rational root, and hence no factor of degree $1$ with rational coefficients.

\textbf{(b)} Subtract and add $x^2$, then group:
\[
x^5+x+1=\left(x^5-x^2\right)+\left(x^2+x+1\right)=x^2(x-1)(x^2+x+1)+(x^2+x+1)=(x^2+x+1)\bigl(x^2(x-1)+1\bigr).
\]
So $\boxed{x^5+x+1=(x^2+x+1)(x^3-x^2+1)}$. Check: $(x^2+x+1)(x^3-x^2+1)=x^5-x^4+x^2+x^4-x^3+x+x^3-x^2+1=x^5+x+1$.

\textbf{Why one would try $x^2+x+1$:} if $\omega$ is a complex cube root of unity, $\omega\ne1$, then $\omega^3=1$ and $\omega^2+\omega+1=0$, so $p(\omega)=\omega^5+\omega+1=\omega^2+\omega+1=0$. Thus $\omega$ and $\bar\omega$ are roots of $p$, and by the factor theorem their quadratic $(x-\omega)(x-\bar\omega)=x^2+x+1$ divides $p(x)$. Finding non-real roots (or clever grouping) can reveal quadratic factors that the rational root theorem cannot.
